% Lecture 9, page 5 of 8 — translation of the 1994 manuscript.
% Fragment intended for appending to courses/course-09.tex.
% Existing sources have not been modified. See editorial note below.

\section*{Lecture 9 - Page 5}

\noindent We denote by $H=\langle a\rangle$ the subgroup of order $n$:
\[
H=\{e,a,\ldots,a^{n-1}\}.
\]

\noindent $b\notin H$.

\noindent Otherwise $b=a^i$. We use the fact that $ba=a^{n-1}b$:
\[
a^{i+1}=a^{n+i-1}\ \big|\ \cdot a^{-i-1}
\quad\Longrightarrow\quad a^{n-2}=e\qquad\times
\]

\bigskip
\noindent $H\cap Hb=\varnothing\ \Longrightarrow\ \underbrace{H\cup Hb}_{2n\ \text{elements}}\subseteq G$.

\bigskip
\noindent We prove that $H\cup Hb=G$, that is,
\[
G=\{e,a,a^2,\ldots,a^{n-1},\ b,ab,a^2b,\ldots,a^{n-1}b\}.
\]

\noindent I show that $H\cup Hb$ is closed under the operation. Since it is finite,
it follows that it is a subgroup.

\[
a^i b\cdot a^s b=a^i a^{s(n-1)}b^2\in\langle a\rangle=H.
\]

\bigskip
\noindent If $x\in G$, then there exist unique $i,j$ such that
\[
0\leq i<n,\qquad 0\leq j<2,
\]
\noindent and $x=a^i b^j$. Let $y=a^s b^t$.\footnote{The exponents in the following second case are overwritten and difficult to decipher in the manuscript. The reading $b^0$ is provisional and has been preserved without mathematical correction, pending confirmation against the original.}

\[
\begin{aligned}
xy&=a^i b^j a^s b^t\\
  &=\begin{cases}
      a^{i+s}b^t, & j=0,\\[2pt]
      a^{i+s(n-1)}b^0, & j=1.
    \end{cases}
\end{aligned}
\]
\bigskip
\noindent $f:G\longrightarrow\Delta_n$,\quad
$f(x)=R^iS^j\ \Longleftrightarrow\ x=a^i b^j$,
\[
0\leq i<n,\qquad 0\leq j<2.
\]
\noindent $f$ is bijective.

\noindent I show that $f(xy)=f(x)f(y)$.

\bigskip
\noindent\underline{\textbf{Action of a group on a set (a $G$-set)}}

\bigskip
\noindent $(G,A,\psi)$,\quad $\psi:G\times A\longrightarrow A$ --- an external
law of composition on $A$ with operators in $G$.

\noindent $G$ --- a group;\quad $A$ --- a non-empty set.

\[
\left\{\begin{aligned}
(1)\quad &\psi(a,\psi(b,u))=\psi(ab,u), && a,b\in G,\ u\in A,\\
(2)\quad &\psi(e,u)=u.
\end{aligned}\right.
\]

\noindent Equivalently, writing $\psi(a,u)=au$:
\[
\left\{\begin{aligned}
(1)\quad &a(bu)=(ab)u,\\
(2)\quad &eu=u.
\end{aligned}\right.
\]

\bigskip
\noindent\underline{Def.}\ Let $G$ be a group and $A$ a set. A map
$\psi:G\times A\longrightarrow A$ is called an \underline{action of the group $G$} on the set $A$
if conditions (1) and (2) are satisfied.
